\section{INTRODUCTION}
Vortex is a unified game library management application with an intelligent recommendation system. The app detects games from multiple platforms into a single interface \& also tracks playtime, session patterns, and user preferences.

Built using C++ and Qt/QML for performance, Vortex implements a machine learning-based recommendation engine that uses content-based filtering and mood-aware personalization. The system analyses user behavior like play duration, time-of-day patterns, and explicit ratings to generate relevant game suggestions.

A modular architecture allows users to customize features according to their needs, avoiding the bloat common in existing launchers. The project demonstrates practical application of recommender systems, data pipeline design, and native software development.

Keywords: Game Library Management, Recommender Systems, Mood-Aware, Personalization, Desktop Application, Machine Learning


\subsection{\underline{Purpose of the Project}}
Digital game distribution has evolved significantly over the past decade. Users now own games across Steam, Epic Games Store, GOG, Xbox Game Pass, and other platforms. Each platform maintains its own launcher, library interface, and usage statistics, creating a fragmented experience for gamers who must switch between multiple applications.

Current game launchers operate in isolation, offering no unified view of a user's complete game collection. Playtime data remains isolated, making it impossible to analyze overall gaming habits. Existing recommendation systems rely primarily on purchase history and popularity metrics rather than actual user behavior and contextual factors like mood or available time.

\subsection{\underline{Target Beneficiary}}
Target users include:

\begin{itemize}
    \item Gamers who own games across multiple digital distribution platforms.
    \item Users who want a centralized game library instead of switching between launchers.
    \item Players interested in tracking gameplay patterns such as playtime and session frequency.
    \item Users who want personalized game recommendations based on mood, time, or previous activity.
\end{itemize}

\subsection{\underline{Project Scope}}
The scope of the Vortex system includes:
\begin{itemize}
    \item Detecting installed games from multiple digital platforms.
    \item Aggregating these games into a unified interface.
    \item Tracking gameplay statistics such as playtime and session history.
    \item Providing a recommendation system based on user behavior.
    \item Offering a modular architecture that allows optional features.
\end{itemize}

The system does not aim to:
\begin{itemize}
    \item Replace existing game launchers.
    \item Modify or distribute game files.
\end{itemize}

\subsection{\underline{References}}
\nocite{*}
\printbibliography

